\section{Introduction}\label{sec:introduction}
Let $k=\C$. An element $f\in k[x_1,\ldots,x_n]$ is a \emph{coordinate} if there are elements $f_2,\ldots,f_n$ such that
\[
 k[x_1,\ldots,x_n]=k[f,f_2,\ldots,f_n]
\]
as a polynomial ring. Equivalently, a polynomial automorphism sends $f$ to $x_1$. It is a \emph{one-stable coordinate} if it is a coordinate in $k[x_1,\ldots,x_n,w]$, where $w$ is independent. The stable coordinate conjecture asks whether an element that becomes a coordinate after adjoining finitely many variables was already a coordinate. The one-variable formulation appears in \cite[Conjecture~3]{ShpilrainYu2002}; successive applications give the formulation with finitely many added variables. Over $\C$, the answer is affirmative in two and three variables by \cite[Proposition~1.1(c)]{ShpilrainYu2002}. The case of one original variable is elementary. Thus four variables are the first possible dimension for a counterexample, and one added variable already suffices.

\begin{theorem}\label{thm:main}
In $R=\C[x_1,x_2,x_3,x_4]$, put
\begin{equation}\label{eq:explicit-counterexample}
 Q=x_2^2-x_4^2+x_1x_3,
 \qquad
 f=x_1-2Q\bigl(Q(x_2+x_4)+x_1x_4\bigr).
\end{equation}
There is a $\C$-algebra automorphism of $R[w]$ sending $f$ to $x_1$, but there is no such automorphism of $R$.
\end{theorem}

The example also has a geometric consequence. A closed embedding $\A^{n-1}\hookrightarrow\A^n$ is \emph{rectifiable} if an automorphism of the ambient affine space takes its image to a coordinate hyperplane. The Abhyankar--Sathaye embedding conjecture asserts, in characteristic zero, that a polynomial $g$ with $k[x_1,\ldots,x_n]/(g)\simeq k^{[n-1]}$ is a coordinate; see \cite[p.~8429]{BlancVanSanten2019}. Here $k^{[r]}$ denotes a polynomial ring in $r$ variables. Corollary~\ref{cor:fibers} gives a negative answer to this conjecture in ambient dimension four. In fact, every fiber is affine three-space and its inclusion in $\A^4$ is nonrectifiable.

\begin{corollary}\label{cor:fibers}
For every $\lambda\in\C$, the hypersurface $f^{-1}(\lambda)\subset\A^4$ is isomorphic to $\A^3$ and its embedding in $\A^4$ is not rectifiable.
\end{corollary}
\begin{proof}
Proposition~\ref{prop:fibers} identifies each fiber with $\A^3$; in particular, $(f-\lambda)$ is its prime defining ideal. If an ambient automorphism takes $f^{-1}(\lambda)$ to a coordinate hyperplane, it takes its defining principal ideal to the ideal of a coordinate. The two generators therefore differ by a nonzero constant. This would make $f-\lambda$, and hence $f$, a coordinate, contrary to Theorem~\ref{thm:main}.
\end{proof}

\subsection*{History and related methods}
The coordinate question retains the embedding in affine space: it asks whether a given function can be made into a coordinate projection by an ambient automorphism. Cancellation asks instead whether an affine variety whose product with an affine line is an affine space must itself be an affine space. Makar-Limanov, van Rossum, Shpilrain, and Yu established an implication from the stable coordinate conjecture to cancellation in \cite[Theorem~1.2]{MakarLimanovRossumShpilrainYu2004}. In our example the ambient variety is already $\A^4$, and the obstruction concerns its distinguished polynomial.

The low-dimensional proof in \cite[proof of Proposition~1.1(c)]{ShpilrainYu2002} uses curve cancellation and the Abhyankar--Moh embedding theorem in two variables. In three variables it combines affine-surface cancellation, using work of Fujita and Miyanishi--Sugie, with Kaliman's theorem that a polynomial on $\C^3$ with general fiber $\C^2$ is a coordinate \cite[Main Theorem]{Kaliman2002}.

The four-variable V\'en\'ereau polynomials have been an important testing ground for these coordinate questions. Lewis proved that the second V\'en\'ereau polynomial is a coordinate \cite{Lewis2013}. Blanc and Poloni developed a description in terms of affine bundles and bivariables, revisited that result, and proved one-stability for a broader family of related polynomials \cite[Proposition~2.3]{BlancPoloni2022}. These results illustrate the distinction between exhibiting polynomial coordinates after stabilization and proving that a polynomial cannot occur in any original coordinate system. For the polynomial in Theorem~\ref{thm:main}, the first problem is settled by explicit substitutions; the second requires a derivation obstruction.

\subsection*{The construction and the obstruction}
The proof uses two presentations of the same four-variable polynomial ring. We first write it as a hypersurface $A=P/(H)$ in a five-variable polynomial ring $P$, with a distinguished element $p$. A determinant-preserving matrix transformation makes $H$ linear in one new coordinate. It identifies $A$ with $R$ and carries $p$ to the polynomial in \eqref{eq:explicit-counterexample}. A second construction, using the exponential of a locally nilpotent derivation, shows that $p$ belongs to polynomial coordinates on $A[w]$. These explicit changes of coordinates occupy Section~\ref{sec:construction}.

To show that $p$ is not a coordinate, we put an integer-indexed filtration on $A$. Its associated graded ring has the form
\[
 G=B[\tau,I/\tau]\subset B[\tau,\tau^{-1}],
 \qquad B=k[x,y,z,u]/(xy-z(z+1)),
\]
where $I$ is an explicitly defined two-generated ideal of $B$ and $\tau$ is the initial form of $p$. The notation $B[\tau,I/\tau]$ means the subalgebra generated by $B$, $\tau$, and all elements $h/\tau$ with $h\in I$. If $p$ were a coordinate, differentiation in a suitable other coordinate would induce a nonzero locally nilpotent derivation $E$ on $B$ whose kernel contains a nonzero element of $I$. Section~\ref{sec:filtration} proves this implication, including the care needed for negative filtration degrees.

The last two sections exclude that invariant. Section~\ref{sec:bundle} passes from $\Spec B$ to the complement of the zero section in a line bundle. In its coordinate ring
\[
 C=k[a,d,b,c,u]/(ac-bd-1),
\]
the ideal $I$ becomes principal, generated by $v=a^3b+a^2u^2-d^2$. The additive action defined by $E$ lifts to this bundle and commutes with fiber scaling. A nonzero invariant in $I$ would force the lifted derivation to annihilate $v$. In Section~\ref{sec:rigidity}, a second grading reduces this to a derivation fixing $a,d,u$. Such a derivation has a coefficient that would simultaneously have to be a polynomial in $a,d,u$ and have negative fiber weight. This is impossible.

The graded-ring and line-bundle strategy adapts the construction in \cite[Sections~2--6]{Cancellation2026}. That paper studies an affine-space cancellation example. Here the quadratic choice in the defining equation has two different uses: it permits the determinant-preserving change that makes $A$ polynomial, and it makes the highest auxiliary part of $v$ factor as $(au-d)(au+d)$. The latter factorization gives the final rigidity argument without a polynomial $abc$ estimate. The passage to leading derivations follows the filtered methods of Kaliman and Makar-Limanov \cite[Sections~3--5]{KalimanMakarLimanov2007}, and the lift of an additive action uses the line-bundle normalization argument of Brion \cite[Lemma~2.9]{Brion2015}. Every construction and obstruction needed for Theorem~\ref{thm:main} is proved below.
